\section{Why Some Markets Make Fortunes Possible}

\subsection{Economic Ceilings and Opportunity Shapes}

Even when 2 founders are equally skilled, their industries may impose very different economic ceilings.

Imagine a neighborhood bakery and a globally distributed software product. The bakery may have devoted customers, excellent quality, and an attractive profit margin. But 1 kitchen, 1 retail location, and a limited radius of customer demand constrain how many orders can be served without substantial investment. The software company may build a product once and sell access to customers in many countries, so long as the software solves a widespread problem and can be reliably supported.

This does not mean software is intrinsically superior to baking. Many software products never find paying users. A strong bakery with multiple locations, a wholesale brand, or packaged consumer products can grow. Software can have high customer acquisition costs, high computing costs, and rapid competitive erosion. What differs is the shape of the opportunity: the potential market, the incremental cost of serving it, the barriers to reaching customers, and the capital required to expand.

\subsection{The Decomposition of Earnings Potential}

A useful decomposition of a firm's earnings potential is:
\begin{equation}
  \Pi = (N \times \text{ARPU}) - C_{\text{operating}} \label{eq:operating-profit}
\end{equation}
where $\Pi$ represents operating profit, $N$ is the number of customers served, $\text{ARPU}$ is the average revenue per customer, and $C_{\text{operating}}$ denotes total operating costs.

This is an accounting identity, not an explanation. The causal questions lie underneath: How many customers exist? How many can we reach? Why will they choose us? How frequently will they buy? How quickly does quality degrade as the customer base grows? How much must we invest to add capacity? What do competitors do when they see our margins?

The answers determine whether a company can grow from 10 customers to 1,000, from 1,000 to 1,000,000, or whether each expansion requires new buildings, staff, permits, and working capital.

\subsection{Capital Intensity and Scalable Capability}

Now introduce profit relative to invested capital. A firm generating ₹1 crore of profit from ₹5 crore of necessary invested capital is different from a firm generating ₹1 crore from ₹50 lakh. The latter may have more freedom to fund growth from internally generated resources. Yet the former may be safer, more durable, or more defensible, depending on the assets involved. Profitability alone does not decide attractiveness; it matters how much capital and risk are required to earn it.

This is 1 reason certain scalable software, licensing, network, and intellectual-property businesses can create extraordinary owner wealth. When a product can reach an enormous market while its incremental cost remains relatively low, the eventual surplus can be large. If customers stay and competitors cannot easily displace the firm, that surplus can persist. But each condition is difficult to achieve. The possibility of spectacular economics invites competition, and success is usually visible only after the fact.

For a founder beginning with little capital, 2 questions are particularly powerful: Can I reach customers without owning expensive infrastructure? And can a portion of what I build for the first customer be reused for the 100th?

A service business might reuse a method, a software component, a reputation, or a specialized team. An education business might reuse course materials, assessments, learning software, and institutional knowledge, while preserving sufficient human attention for actual learning. A consumer brand might reuse packaging design, supplier relationships, distribution agreements, and repeat purchasing habits. \textbf{What scales is not always the product itself. Sometimes it is the capability to deliver the product.}
